\documentclass[11pt,a4paper]{article}
\usepackage[utf8]{inputenc}
\usepackage[T1]{fontenc}
\usepackage[brazilian]{babel}
\usepackage{geometry}
\usepackage{graphicx}
\usepackage{booktabs}
\usepackage{float}
\usepackage{enumitem}
\usepackage{amsmath}
\usepackage{xcolor}
\usepackage{tcolorbox}
\geometry{margin=2.5cm}
\setlength{\parindent}{0pt}
\setlength{\parskip}{0.5em}
\tcbuselibrary{skins,breakable}
\newtcolorbox{nota}[1][]{colback=yellow!10,colframe=orange!60!black,fonttitle=\bfseries,title=#1,breakable}

\begin{document}
\begin{center}
  
\includegraphics[width=0.16\textwidth]{../../../figures_global/unb.png}\\[0.3cm]
  {\large\textbf{GABARITO DO PROFESSOR -- USO RESTRITO}}\\[0.2cm]
  {\large FGA0092 -- Princípios de Comunicação para Engenharia}\\[0.2cm]
  \rule{12cm}{0.8pt}\\[0.2cm]
  {\Large\textbf{Prática 07: Phase-Locked Loop (PLL)}}\\[0.2cm]
  \rule{12cm}{0.8pt}
\end{center}

\section*{Gráfico esperado}
\begin{center}
  \includegraphics[width=0.88\textwidth]{grafico_esperado.png}
\end{center}
\textit{O gráfico acima foi gerado pelo script \texttt{gerar\_grafico.py} desta pasta.}

\section*{Tabelas preenchidas (valores de referência)}

Os valores abaixo foram obtidos pela simulação headless em
\texttt{Laboratorio/Codigos/pratica\_06\_gabarito.py}
($f_s = 32000$\,Hz, janela de 1\,s, \texttt{Max/Min Phase/sample} =
$\pm 2\pi \cdot 2000/f_s$). A coluna \emph{Locked} usa o critério
$\sigma/|\mu| < 0{,}02$ na cauda do sinal.

\begin{table}[H]
  \centering
  \begin{tabular}{c c c c}
    \toprule
    $f_\text{in}$ (Hz) & Valor de regime (rad/amostra) & $f_\text{est}$ (Hz) & Locked \\
    \midrule
    500  & 0{,}098175 & 500{,}00  & 1 \\
    800  & 0{,}157080 & 800{,}00  & 1 \\
    1000 & 0{,}196350 & 1000{,}00 & 1 \\
    1200 & 0{,}235619 & 1200{,}00 & 1 \\
    1500 & 0{,}294524 & 1500{,}00 & 1 \\
    \bottomrule
  \end{tabular}
  \caption{Tabela~1 (\texttt{tab:lock}) -- saída do PLL em função de $f_\text{in}$ com \texttt{loop\_bw} = 0{,}062. Valor esperado em rad/amostra = $2\pi f_\text{in}/f_s$; $f_\text{est} = \mu \cdot f_s/(2\pi)$ deve reproduzir $f_\text{in}$ exatamente quando há lock.}
\end{table}

\begin{table}[H]
  \centering
  \begin{tabular}{c c c}
    \toprule
    \texttt{loop\_bw} & $t_\text{conv}$ (ms) & Overshoot (\%) \\
    \midrule
    0{,}010 & 18{,}12 & 2{,}5 \\
    0{,}062 & 1{,}53  & 4{,}4 \\
    0{,}200 & 0{,}50  & 4{,}7 \\
    0{,}500 & 0{,}22  & 7{,}4 \\
    \bottomrule
  \end{tabular}
  \caption{Tabela~2 (\texttt{tab:pll}) -- métricas em função de \texttt{loop\_bw}, com $f_\text{in} = 1000$\,Hz.}
\end{table}

\begin{table}[H]
  \centering
  \begin{tabular}{c c c c}
    \toprule
    $f_\text{in}$ (Hz) & Valor de regime (rad/amostra) & $f_\text{est}$ (Hz) & Locked \\
    \midrule
    200  & 0{,}039270 & 200{,}00  & 1 \\
    600  & 0{,}117810 & 600{,}00  & 1 \\
    1000 & 0{,}196350 & 1000{,}00 & 1 \\
    1400 & 0{,}274889 & 1400{,}00 & 1 \\
    1800 & 0{,}353429 & 1800{,}00 & 1 \\
    \midrule
    \multicolumn{3}{l}{$\Delta f_\text{cap}$ medida (Hz)} & 1600 \\
    \multicolumn{3}{l}{$\Delta f_\text{cap}$ teórica (Hz)} & 631{,}5 \\
    \bottomrule
  \end{tabular}
  \caption{Tabela~3 (\texttt{tab:capture}) -- faixa de captura com \texttt{loop\_bw} = 0{,}062. $f_\text{est}$ reproduz $f_\text{in}$ em todas as cinco linhas, confirmando lock. A medida supera a teórica porque a janela de 2\,s permite \emph{pull-in} além da faixa estrita de captura; o limite duro vem de \texttt{Max/Min Phase/sample} = $\pm 2000$\,Hz.}
\end{table}

\begin{table}[H]
  \centering
  \begin{tabular}{c c c}
    \toprule
    \texttt{noise\_amp} & Jitter std \texttt{loop\_bw}=0{,}062 & Jitter std \texttt{loop\_bw}=0{,}200 \\
    \midrule
    0{,}0 & 0{,}000000 & 0{,}000000 \\
    0{,}1 & 0{,}001875 & 0{,}010717 \\
    0{,}3 & 0{,}005685 & 0{,}032819 \\
    0{,}5 & 0{,}010033 & 0{,}058452 \\
    1{,}0 & 0{,}023289 & 0{,}140502 \\
    \bottomrule
  \end{tabular}
  \caption{Tabela~4 (\texttt{tab:ruido}) -- jitter de fase em função de \texttt{noise\_amp}, com $f_\text{in} = 1000$\,Hz. Loop largo (0{,}200) gera $\sim$5,8x mais jitter que o loop estreito (0{,}062).}
\end{table}

\begin{nota}[Tolerâncias de correção]
Aceitar variações de até $\pm 5\%$ nos valores de regime, $\pm 20\%$
no tempo de convergência, $\pm 30\%$ no overshoot e $\pm 20\%$ no jitter
(o \texttt{Noise Source} depende da \texttt{Seed}).
\end{nota}

\section*{Tabelas das análises complementares}

\begin{table}[H]
  \centering
  \begin{tabular}{c c c}
    \toprule
    \texttt{noise\_amp} & $\sigma_f$ \texttt{lbw}=0{,}062 (Hz) & $\sigma_f$ \texttt{lbw}=0{,}200 (Hz) \\
    \midrule
    0{,}0 & 0{,}00   & 0{,}00   \\
    0{,}1 & 9{,}55   & 54{,}58  \\
    0{,}3 & 28{,}95  & 167{,}15 \\
    0{,}5 & 51{,}10  & 297{,}69 \\
    1{,}0 & 118{,}61 & 715{,}57 \\
    \bottomrule
  \end{tabular}
  \caption{Tabela~5 (\texttt{tab:jitter\_hz}) -- jitter da estimativa em Hz. Os mesmos números da Tabela~4 multiplicados por $f_s/(2\pi)$. Com loop largo e ruído alto, a incerteza na estimativa de frequência ultrapassa 700~Hz -- maior que a própria $f_\text{in}$.}
\end{table}

\begin{table}[H]
  \centering
  \begin{tabular}{c c c c}
    \toprule
    $f_\text{in}$ (Hz) & $f_\text{est}$ (Hz) & Desvio (Hz) & Locked \\
    \midrule
    1500 & 1500{,}00 & $+0{,}00$ & 1 \\
    1900 & 1900{,}00 & $+0{,}00$ & 1 \\
    2000 & 2000{,}00 & $-0{,}00$ & 1 \\
    2100 & 2000{,}00 & $-100{,}00$ & 1 \\
    2300 & 2000{,}00 & $-300{,}00$ & 1 \\
    \bottomrule
  \end{tabular}
  \caption{Tabela~6 (\texttt{tab:saturacao}) -- saturação no limite de hold. Para $f_\text{in} > 2000$\,Hz o $f_\text{est}$ trava em 2000\,Hz porque o bloco \texttt{PLL Frequency Detector} foi configurado com \texttt{Max Phase/sample} = $2\pi \cdot 2000/f_s$. Observe que \emph{locked=1} mesmo nas linhas saturadas: o critério $\sigma/|\mu| < 0{,}02$ não distingue lock real de saturação no limite, porque ambos produzem saída constante. Discuta isso com o aluno.}
\end{table}

\begin{table}[H]
  \centering
  \begin{tabular}{c c c}
    \toprule
    $f_\text{in}$ (Hz) & $f_\text{est}$ janela 50 ms (Hz) & Locked \\
    \midrule
    200  & 200{,}00  & 1 \\
    600  & 600{,}00  & 1 \\
    1000 & 1000{,}00 & 1 \\
    1400 & 1400{,}00 & 1 \\
    1800 & 1800{,}00 & 1 \\
    \midrule
    \multicolumn{2}{l}{$\Delta f_\text{cap}$ janela curta (Hz)} & 1600 \\
    \multicolumn{2}{l}{$\Delta f_\text{cap}$ teórica (Hz)} & 631{,}5 \\
    \bottomrule
  \end{tabular}
  \caption{Tabela~7 (\texttt{tab:janela\_curta}) -- captura em 50\,ms. Resultado pedagógico: a faixa medida continua em 1600\,Hz, igual à janela de 2\,s, mostrando que a implementação \texttt{analog.pll\_freqdet\_cf} captura \emph{toda} a faixa de hold em menos de 50\,ms. A previsão teórica $2\omega_n/2\pi \approx 631$\,Hz é conservadora: vale para PLLs lineares de 2ª ordem, mas a implementação no GNU Radio usa um \texttt{control\_loop} com alpha e beta ajustados que faz captura agressiva. Para reproduzir o regime teórico seria preciso reduzir a janela para algo da ordem de $t_\text{conv}$ ($\sim$ 1,5\,ms) -- abaixo disso a estatística fica pobre demais.}
\end{table}

\begin{table}[H]
  \centering
  \begin{tabular}{c c c}
    \toprule
    \texttt{noise\_amp} & $f_\text{est}$ \texttt{lbw}=0{,}062 (Hz) & $f_\text{est}$ \texttt{lbw}=0{,}200 (Hz) \\
    \midrule
    0{,}0 & 1000{,}00 & 1000{,}00 \\
    0{,}3 & 1000{,}08 & 999{,}07  \\
    0{,}5 & 999{,}92  & 998{,}63  \\
    1{,}0 & 1001{,}06 & 911{,}39  \\
    \bottomrule
  \end{tabular}
  \caption{Tabela~8 (\texttt{tab:bias\_curto}) -- $f_\text{est}$ em janela curta de 50\,ms com $f_\text{in} = 1000$\,Hz. Sem ruído, $f_\text{est}$ é exato. Com \texttt{noise\_amp} = 1{,}0 e loop largo, o bias chega a $\sim$89\,Hz -- aproximadamente uma década pior que o loop estreito. Mostra que com pouco tempo de observação a média não tem amostras suficientes para "lavar" o jitter, e a estimativa de uma única medida vira pouco confiável. Em telecom isso é a base da escolha de constante de tempo do PLL em sistemas de aquisição de portadora.}
\end{table}

\begin{nota}[Tolerâncias das tabelas complementares]
Tabela~5 deriva exatamente da~4, então mesmas tolerâncias. Tabela~6
deve ser idêntica para qualquer aluno (sem ruído). Tabelas~7 e~8 podem
variar até $\pm 50\,\%$ por causa da \texttt{Seed} do ruído e da
sensibilidade da janela curta -- avaliar o raciocínio, não os números
exatos.
\end{nota}

\section*{Respostas esperadas}

\begin{enumerate}[label=\textbf{Q\arabic*.}, leftmargin=*, itemsep=1em]

  \item \textbf{Descreva os três componentes fundamentais de um PLL e a função de cada um.}

    \textbf{Resposta esperada:} (1) Detector de Fase (PD): multiplica o sinal de entrada
    pelo sinal do VCO e filtra, produzindo um sinal proporcional a
    $\sin(\phi_i - \phi_o) \approx \phi_e$ para pequenos erros de fase.
    (2) Filtro de Laço (LF): filtro passa-baixas que remove altas frequências e
    determina a largura de banda do sistema (velocidade vs.\ supressão de ruído).
    (3) VCO: oscilador com frequência $\omega_0 + K_v \cdot e(t)$; integra o sinal
    de erro para ajustar a fase. A realimentação fecha o laço: PD detecta erro, LF
    suaviza, VCO corrige.

    \begin{nota}[Erros comuns]
      Descrever apenas o VCO e o comparador de fase sem mencionar o papel do filtro
      de laço na determinação da largura de banda e na supressão de ruído.
    \end{nota}

  \item \textbf{Como se detecta o estado de lock do PLL na saída do \texttt{PLL Freq Det}?}

    \textbf{Resposta esperada:} O PLL está em lock quando a saída do \texttt{PLL Freq Det}
    converge para um nível DC proporcional a $\omega_\text{in} = 2\pi f_\text{in}/f_s$.
    Critério numérico aplicado no pós-processamento: pegar os 20\,\% finais do vetor de
    saída e exigir $\sigma/|\mu| < 0{,}02$, ou seja, ruído residual menor que 2\,\% do
    valor médio. Na Tabela~1, todas as cinco frequências satisfazem esse critério: o
    valor de regime é exatamente $2\pi f_\text{in}/f_s$, por exemplo $f_\text{in}=1000$
    Hz dá $\mu = 0{,}196350$ rad/amostra, idêntico a $2\pi \cdot 1000/32000$.

    \begin{nota}[Erros comuns]
      Confundir a convergência da saída com o desaparecimento completo de oscilações
      transitórias. O sobresinal pode persistir brevemente mesmo em lock. Outro erro
      é declarar lock só observando o Time Sink a olho, sem o critério numérico aplicado
      à cauda do sinal salvo via File Sink.
    \end{nota}

  \item \textbf{Como a largura de banda do laço (\texttt{loop\_bw}) afeta a velocidade de lock e o jitter de fase?}

    \textbf{Resposta esperada:} A Tabela~2 mostra o tradeoff clássico: aumentando
    \texttt{loop\_bw} de 0{,}010 para 0{,}500, $t_\text{conv}$ cai de 18{,}12 ms para
    0{,}22 ms -- duas ordens de grandeza mais rápido -- enquanto o overshoot sobe
    de 2{,}5\,\% para 7{,}4\,\%, indicando resposta cada vez mais subamortecida.
    O valor padrão $0{,}062$\,rad/amostra fica próximo do \emph{joelho} do
    tradeoff: convergência em 1{,}53 ms com overshoot moderado de 4{,}4\,\%.
    Sem ruído na entrada, o efeito sobre o jitter pós-lock não é mensurável
    nesta etapa -- isso aparece quantitativamente na Etapa~5.

    \begin{nota}[Erros comuns]
      Afirmar que \texttt{loop\_bw} larga é sempre preferível por convergir mais rápido,
      sem mencionar o aumento de jitter de fase e a possível instabilidade. Ou
      ignorar a Tabela~2 e responder qualitativamente sem citar valores.
    \end{nota}

  \item \textbf{Qual foi a faixa de captura medida e como ela se compara à previsão teórica?}

    \textbf{Resposta esperada:} Da Tabela~3, todas as cinco frequências de 200 a 1800
    Hz travaram, então $\Delta f_\text{cap,\,medida} = 1800 - 200 = 1600$\,Hz. A
    previsão teórica é $\Delta f_\text{cap} \approx \texttt{loop\_bw}\cdot f_s/\pi
    = 0{,}062 \cdot 32000/\pi \approx 631{,}5$\,Hz. A medida é cerca de
    $2{,}5\times$ maior que a teórica e isso \emph{não} é erro: o limite teórico
    descreve a faixa de captura rápida -- pull-in instantâneo no segundo regime --,
    mas a janela de 2\,s usada na simulação dá ao laço tempo suficiente para fazer
    \emph{pull-in lento}, que estende o range efetivo até o limite duro imposto
    por \texttt{Max/Min Phase/sample} = $\pm 2000$\,Hz. Discrepâncias adicionais
    vêm da não-linearidade do detector de fase ($\sin(\phi)$ em vez de $\phi$) e
    da quantização discreta do filtro de laço.

    \begin{nota}[Erros comuns]
      Calcular a faixa de captura em rad/amostra sem converter para Hz, ou usar
      \texttt{loop\_bw} diretamente como a faixa de captura. Ou tratar a
      discrepância de $2{,}5\times$ como erro experimental em vez de
      reconhecer o pull-in lento.
    \end{nota}

  \item \textbf{Como o ruído externo afeta o jitter de fase do PLL em função de \texttt{loop\_bw}?}

    \textbf{Resposta esperada:} A Tabela~4 mostra que o jitter cresce linearmente
    com \texttt{noise\_amp} e que, para o mesmo \texttt{noise\_amp}, o laço largo
    sempre produz mais jitter que o estreito: por exemplo em \texttt{noise\_amp}=0{,}5
    o jitter é 0{,}010033 com \texttt{loop\_bw}=0{,}062 e 0{,}058452 com
    \texttt{loop\_bw}=0{,}200, ou seja, $\sim 5{,}8\times$ maior. Isso é direto da
    teoria: a variância de saída do PLL escala como $\sigma^2_{\phi_o} \propto
    \sigma^2_n \cdot B_L$, onde $B_L$ é a banda equivalente de ruído do laço.
    Laço largo deixa passar mais conteúdo espectral do ruído de entrada; laço
    estreito atua como filtro passa-baixas mais agressivo. O tradeoff com a
    velocidade de aquisição limita o quão estreito é viável usar.

    \begin{nota}[Erros comuns]
      Afirmar que \texttt{loop\_bw} estreita é sempre melhor para ruído, sem mencionar
      o tradeoff com a velocidade de aquisição de lock e o risco de perda de lock com
      ruído intenso. Ou apresentar a Tabela~4 sem extrair a razão $\sim 5{,}8\times$
      entre os dois \texttt{loop\_bw}.
    \end{nota}

\end{enumerate}

\section*{Explicação esperada do pós-processamento}

O entregável exige que o aluno explique em palavras próprias o que cada
snippet de Python está fazendo. As respostas mínimas aceitáveis são:

\textbf{Etapa 2 -- Tabela~1 (lock).} \texttt{np.fromfile} interpreta o
\texttt{.bin} como vetor de \texttt{float32}, mesmo formato escrito pelo
\texttt{File Sink} configurado como Float. O slice \texttt{data[int(0.8*len(data)):]}
descarta o transitório inicial e mantém os 20\,\% finais, onde o PLL já está em
regime. \texttt{np.mean(tail)} é o valor de regime $\mu$ e \texttt{np.std(tail)}
é a flutuação $\sigma$ em torno desse valor. O critério $\sigma/|\mu| < 0{,}02$
declara lock quando o jitter relativo está abaixo de 2\,\% -- o caso especial
$|\mu| < 10^{-6}$ trata o limite numérico em que o PLL não converge.

\textbf{Etapa 3 -- Tabela~2 (loop\_bw).} O valor final é estimado pelos 10\,\%
finais, dando uma referência mais limpa. A máscara booleana \texttt{within} marca
amostras dentro de uma faixa de 5\,\% do valor final. O \texttt{next(...)} encontra
o primeiro índice em que essa condição se sustenta por 100 amostras consecutivas
-- evita gatilhos prematuros nas oscilações do transitório -- e o índice é
convertido em milissegundos por $t/f_s$. O overshoot é
$(\max|data| - |\mu|)/|\mu| \cdot 100\,\%$, capturando o quanto o pico
ultrapassou o valor final.

\textbf{Etapa 4 -- Tabela~3 (capture).} O mesmo critério da Etapa 2 é aplicado
por frequência. A lista \texttt{locked\_freqs} acumula apenas as freqs com lock;
$\Delta f_\text{cap} = \max(\text{locked}) - \min(\text{locked})$. A teórica é
\texttt{loop\_bw} $\cdot f_s/\pi$, derivada de
$\Delta\omega_\text{cap} \approx 2\omega_n$.

\textbf{Etapa 5 -- Tabela~4 (ruído).} Loop duplo sobre os pares
$(\texttt{noise\_amp}, \texttt{loop\_bw})$. Só o desvio padrão da cauda é
necessário -- o valor médio é fixo em $2\pi \cdot 1000/32000$ porque
$f_\text{in}$ é fixo. O jitter cresce com \texttt{noise\_amp} e com
\texttt{loop\_bw}: o aluno deve articular que ambos efeitos vêm da mesma causa,
o filtro de laço admitir mais conteúdo espectral do ruído de entrada.

\begin{nota}[Erros comuns na explicação]
Limitar-se a parafrasear nomes de funções -- ``\texttt{np.mean} calcula a média''
-- sem dizer \emph{o que} essa média representa fisicamente. Ignorar
\emph{por que} os 20\,\% finais são usados, ou trocar o critério
$\sigma/|\mu|$ por inspeção visual.
\end{nota}

\section*{Critérios de correção}
\begin{tabular}{lc}
  \toprule
  Critério & Pontuação \\
  \midrule
  Print do flowgraph completo, incluindo o \texttt{File Sink}    & 0,5 ponto \\
  Tabela~1 (lock) preenchida e consistente com $2\pi f_\text{in}/f_s$ & 0,5 ponto \\
  Tabela~2 (loop\_bw) preenchida com $t_\text{conv}$ e overshoot & 0,5 ponto \\
  Tabela~3 (capture) preenchida + $\Delta f_\text{cap}$ medida/teórica & 0,5 ponto \\
  Tabela~4 (ruído) preenchida nos dois \texttt{loop\_bw}          & 0,5 ponto \\
  Explicação coerente do pós-processamento das Etapas 2--5         & 1,0 ponto \\
  Q1 -- Três componentes com funções corretas                       & 0,3 ponto \\
  Q2 -- Critério de lock numérico ($\sigma/|\mu|$) bem articulado   & 0,3 ponto \\
  Q3 -- Tradeoff velocidade vs.\ ruído citando valores da Tabela~2  & 0,3 ponto \\
  Q4 -- Comparação medida vs.\ teórica + discussão do pull-in lento & 0,3 ponto \\
  Q5 -- Relação jitter vs.\ \texttt{loop\_bw} citando a Tabela~4    & 0,3 ponto \\
  \midrule
  \textbf{Total} & \textbf{5,0 pontos} \\
  \bottomrule
\end{tabular}

\bigskip
\begin{nota}[Nota de correção geral]
  Priorizar coerência física e evidência numérica das tabelas. Aceitar
  variações de até $\pm 5\,\%$ nos valores de regime, $\pm 20\,\%$ no tempo
  de convergência, $\pm 30\,\%$ no overshoot e $\pm 20\,\%$ no jitter -- a
  \texttt{Seed} do \texttt{Noise Source} influencia o jitter. Desconto duplo
  para tabelas preenchidas sem que o aluno explique o código de pós-processamento.
\end{nota}

\end{document}
